\section{Сведение к задаче целочисленного линейного программирования}
\label{sec:4_reduce_description} \index{4_reduce_description}

\subsection*{Входные данные}
Каждая работа $i$ ($i=1,\dots,n$) характеризуется длительностью $d_i>0$ и набором буферов, которые она создаёт:
\[
V = \bigl\{(i,\; n_i,\; r_{i1}, r_{i2}, \dots, r_{i n_i})\bigr\}_{i=1}^{n},
\]
где $n_i$ — количество буферов работы $i$, а $r_{ij}>0$ — объём памяти, занимаемый $j$-м буфером.  
Дуги графа задают передачу данных:
\[
E = \{(i, k, j)\},\quad j\in[1,n_i],
\]
т.е. дуга $(i,k,j)$ означает, что $j$-й буфер работы $i$ используется работой $k$ (является её входным буфером).  
Все времена $d_i$ целые, допустимый объём памяти равен $M$, число процессоров $P$.

\subsection*{Переменные}

\begin{itemize}
	\item $s_i$ — время старта работы $i$ (целочисленная переменная).
	\item $m_{ij}$ — бинарная переменная, $m_{ij}=1$ если $s_i \le s_j$, иначе $0$ ($i,j=1,\dots,n$).
	\item $p_{ri}$ — бинарная переменная, $p_{ri}=1$ если работа $i$ назначена на процессор $r$ ($r=1,\dots,P$).
	\item $t_{ij}$ — бинарная переменная, $t_{ij}=1$ если работы $i$ и $j$ выполняются на одном процессоре.
	\item $l$ — целочисленная переменная, общая длительность расписания (makespan).
	\item Для каждого буфера $(i,j)$ ($j=1,\dots,n_i$) вводятся целочисленные переменные:
	\[
	\text{start}_{ij},\quad \text{max\_end}_{ij}.
	\]
	$\text{start}_{ij}$ совпадает со стартом работы $i$, а $\text{max\_end}_{ij}$ — максимальное время завершения среди всех потребителей этого буфера.
	\item Для каждой пары буферов $(i,j)$ и $(p,q)$ вводятся бинарные переменные:
	\[
	z_{ij}^{pq},\; zz_{ij}^{pq},\; a_{ij}^{pq},\; aa_{ij}^{pq}.
	\]
	Они служат для моделирования пересечений интервалов активности буферов.
\end{itemize}

\subsection*{Вспомогательная константа}

\[
D = \sum_{i=1}^{n} d_i,
\]
которая заведомо больше любого возможного makespan.

\subsection*{Ограничения, связывающие $m_{ij}$ и $s_i$}

Для всех $i,j$:
\begin{align}
	m_{ii} &= 1, \label{eq:m_diag}\\
	D\,m_{ij} - s_j + s_i &\ge 0, \label{eq:m_ge}\\
	s_j - s_i - D\,m_{ij} &\ge 1 - D. \label{eq:m_le}
\end{align}
Эти условия эквивалентны: $m_{ij}=1 \Leftrightarrow s_i \le s_j$.

\subsection*{Назначение на процессоры}

Для всех $i$ и $r$:
\begin{align}
	\sum_{r=1}^{P} p_{ri} &= 1, \label{eq:p_sum}\\
	p_{ri} &\in \{0,1\}. \nonumber
\end{align}

\subsection*{Переменные $t_{ij}$ (один процессор)}

Для всех $i,j,k$ ($i\neq j$) и всех $r$:
\begin{align}
	t_{ij} &\ge p_{ri} + p_{rj} - 1, \label{eq:t_ge}\\
	t_{ij} &\le p_{ri} - p_{rj} + 1, \label{eq:t_le}\\
	t_{ij} &\in \{0,1\}. \nonumber
\end{align}
При этом $t_{ij}=1$ тогда и только тогда, когда $p_{ri}=p_{rj}=1$ для некоторого $r$.

\subsection*{Зависимости по графу}

Для каждой дуги $(i,k,j)\in E$ (т.е. $j$-й буфер работы $i$ потребляется работой $k$):
\begin{align}
	s_k \ge s_i + d_i. \label{eq:prec}
\end{align}

\subsection*{Определение $\text{start}_{ij}$ и $\text{max\_end}_{ij}$}

Для каждого буфера $(i,j)$:
\begin{align}
	\text{start}_{ij} &= s_i, \label{eq:start_def}\\
	\text{max\_end}_{ij} &\ge s_k + d_k \quad \text{для всех } k \text{ таких, что } (i,k,j)\in E. \label{eq:max_end_def}
\end{align}
Таким образом, $\text{max\_end}_{ij}$ — максимум моментов завершения всех потребителей буфера.

\subsection*{Взаимное исключение на одном процессоре}

Для всех $i\neq j$:
\begin{align}
	s_i - s_j + D\,m_{ij} - D\,t_{ij} &\ge -D + d_j, \label{eq:no_conflict1}\\
	s_j - s_i - D\,m_{ij} - D\,t_{ij} &\ge -2D + d_i. \label{eq:no_conflict2}
\end{align}
Если $t_{ij}=1$ (работы на одном процессоре), то эти ограничения вынуждают $s_j \ge s_i + d_i$ при $m_{ij}=1$ и $s_i \ge s_j + d_j$ при $m_{ij}=0$. При $t_{ij}=0$ ограничения выполняются автоматически.

\subsection*{Учёт пересечений буферов (для ограничения памяти)}

Для каждой пары буферов $(i,j)$ и $(p,q)$ введём вспомогательные бинарные переменные $z_{ij}^{pq}$, $zz_{ij}^{pq}$, $a_{ij}^{pq}$, $aa_{ij}^{pq}$, связанные следующими линейными ограничениями:
\begin{align}
	\text{start}_{pq} - \text{start}_{ij} - (D+1)\,z_{ij}^{pq} &\ge -D, \label{eq:z_start}\\
	(D+1)\,a_{ij}^{pq} + (D+1)\,z_{ij}^{pq} + \text{start}_{ij} - \text{max\_end}_{pq} &\ge 0, \label{eq:a_start}\\
	\text{start}_{pq} - \text{max\_end}_{ij} - (D+1)\,zz_{ij}^{pq} &\ge -D, \label{eq:zz_end}\\
	(D+1)\,aa_{ij}^{pq} + (D+1)\,zz_{ij}^{pq} + \text{max\_end}_{ij} - \text{max\_end}_{pq} &\ge 0. \label{eq:aa_end}
\end{align}
Переменные $a_{ij}^{pq}$ и $aa_{ij}^{pq}$ интерпретируются следующим образом:
\begin{itemize}
	\item $a_{ij}^{pq}=1$, если буфер $(p,q)$ активен в момент старта буфера $(i,j)$ (т.е. $\text{start}_{ij}$ лежит внутри интервала активности $(p,q)$);
	\item $aa_{ij}^{pq}=1$, если буфер $(p,q)$ активен в момент завершения буфера $(i,j)$ (т.е. $\text{max\_end}_{ij}$ лежит внутри интервала активности $(p,q)$).
\end{itemize}
Интервал активности буфера $(p,q)$ — это $[\text{start}_{pq},\; \text{max\_end}_{pq})$.

\subsection*{Ограничение на суммарную память}

Для каждого буфера $(i,j)$ суммарный объём памяти, занятый в момент $\text{start}_{ij}$ (соответственно $\text{max\_end}_{ij}$), не должен превышать $M$:
\begin{align}
	\sum_{(p,q)} r_{pq}\, a_{ij}^{pq} &\le M, \label{eq:mem_start}\\
	\sum_{(p,q)} r_{pq}\, aa_{ij}^{pq} &\le M, \label{eq:mem_end}
\end{align}
где сумма берётся по всем буферам, присутствующим в графе. Благодаря определению $a_{ij}^{pq}$ и $aa_{ij}^{pq}$, эти ограничения гарантируют, что пиковое потребление памяти в любой момент (все точки вида $\text{start}_{ij}$ или $\text{max\_end}_{ij}$) не превышает $M$.

\subsection*{Целевая функция}

Минимизируем общую длительность расписания:
\begin{align}
	\min \; l,\\
	l \ge s_i + d_i \quad \forall i. \label{eq:l_bound}
\end{align}

\subsection*{Итоговая задача}

\[
\begin{array}{ll}
	\text{Минимизировать} & l \\
	\text{при ограничениях} & (\ref{eq:m_diag})—(\ref{eq:m_le}),\; (\ref{eq:p_sum}),\; (\ref{eq:t_ge})—(\ref{eq:t_le}),\; (\ref{eq:prec}),\; (\ref{eq:start_def}),\; (\ref{eq:max_end_def}),\\
	& (\ref{eq:no_conflict1})—(\ref{eq:no_conflict2}),\; (\ref{eq:z_start})—(\ref{eq:aa_end}),\; (\ref{eq:mem_start})—(\ref{eq:mem_end}),\; (\ref{eq:l_bound}),\\
	& \text{переменные } s_i,\text{start}_{ij},\text{max\_end}_{ij},l\in\mathbb{Z}_{\ge0},\; 
	m_{ij},p_{ri},t_{ij},z_{ij}^{pq},zz_{ij}^{pq},a_{ij}^{pq},aa_{ij}^{pq}\in\{0,1\}.
\end{array}
\]

\subsection*{Число переменных и ограничений}

Пусть $n$ — число работ, $n_i$ — число буферов работы $i$, $B=\sum_i n_i$ — общее число буферов, $P$ — число процессоров. Тогда:
\begin{itemize}
	\item $m_{ij}$: $n^2$ бинарных.
	\item $s_i$: $n$ целых.
	\item $p_{ri}$: $P n$ бинарных.
	\item $t_{ij}$: $n^2$ бинарных.
	\item $\text{start}_{ij},\text{max\_end}_{ij}$: $2B$ целых.
	\item $z_{ij}^{pq}, zz_{ij}^{pq}, a_{ij}^{pq}, aa_{ij}^{pq}$: $4B^2$ бинарных.
	\item $l$: одна целая.
\end{itemize}
Количество ограничений составляет $O(n^2 + P n + |E| + B^2)$.